\begin{abstract}
Center balancing and multiple class prototypes are simple choices for transferring frozen visual representations across institutions. We study whether they improve recognition when features, partitions, and source validation are controlled. Our study implementation, center-balanced multi-prototype classification (CBMP), combines inverse group-size weighting, class-conditional center balancing, and spherical clustering. Eight controls separate these choices from discriminative linear and cosine fitting. This design isolates classifier choices within each encoder configuration. We evaluate reviewed PolypGen sequences with ResNet18 and DINOv2 and a custom CAMELYON17-WILDS cohort with ResNet18, UNI, and CONCH. Each comparison uses ten fitting seeds, historical grids, and a four-candidate sensitivity protocol. Under four-candidate selection, primary-hospital trained cosine reaches 97.96\% balanced accuracy with UNI and 96.88\% with CONCH, compared with CBMP scores of 95.80\% and 92.63\%. Source validation selects exactly one center-balanced prototype per class for CONCH in all ten seeds, whereas UNI selects four prototypes. Stronger histology representations therefore do not establish a consistent advantage for prototype balancing or multiplicity. Group-pair audits distinguish within-group discrimination from across-group score comparability; patient-macro summaries expose dependence on evaluation weighting. The findings remain conditional on observed cohorts. Full official benchmark evaluation, auditing of pretraining overlap, and independent patient-linked confirmation are required before broader institutional robustness claims.
\end{abstract}

